%% ============================================================
%%  Chapter 1 — Introduction
%% ============================================================
\chapter{Introduction}
\label{ch:introduction}

\section{Background and Motivation}
\label{sec:motivation}

The rapid adoption of cloud-native architectures and container orchestration platforms,
particularly Kubernetes~\cite{kubernetes2014}, has fundamentally transformed how modern
applications are designed, deployed, and operated. Microservice architectures, which
decompose monolithic applications into loosely coupled, independently deployable services,
have become the dominant paradigm for large-scale distributed systems~\cite{microservices_survey}.

This shift, however, has introduced a complex and evolving threat landscape. Distributed
Denial of Service (DDoS) attacks remain one of the most pervasive and economically
damaging classes of network attacks~\cite{cloudflare_threat_report_2024}. In cloud-native
environments, the challenge is compounded because attack traffic can originate from within
the cluster itself (compromised pods), traverse encrypted service mesh tunnels, and exploit
the dynamic, ephemeral nature of containerised workloads.

Existing DDoS defense mechanisms operate predominantly at the network edge ---
inspecting individual packets, computing statistical flow features, or applying
signature-based filters~\cite{aws_shield, cloudflare_magic_transit}. While effective
against volumetric floods, these approaches share a fundamental blind spot: they
\emph{treat traffic as a collection of independent flows} rather than as the product of an
interconnected system of communicating services. Consequently, they fail to model the
relational and temporal structure of microservice communication, limiting their ability to
detect sophisticated, low-and-slow, and application-layer attacks that mimic legitimate
traffic behaviour.

Recent advances in \emph{eBPF (Extended Berkeley Packet Filter)}~\cite{gregg2019bpf}
have enabled unprecedented observability into kernel-level network events with near-zero
overhead, providing a foundation for fine-grained, in-cluster telemetry. Simultaneously,
\emph{Graph Neural Networks (GNNs)}~\cite{zhou2020gnn_survey} have emerged as a powerful
paradigm for learning on relational data, demonstrating strong performance in
fraud detection, social network analysis, and network intrusion detection tasks.

This project proposes to bridge these two advances: leveraging eBPF-based telemetry from
within a Kubernetes cluster to construct \emph{dynamic communication graphs} of
microservice interactions, and applying graph-based anomaly detection to identify DDoS
attack patterns that are invisible to traditional, flow-centric approaches.

% ─────────────────────────────────────────────
\section{Problem Statement}
\label{sec:problem}

The core problem addressed by this research is:

\begin{quote}
\textit{How can the relational structure of Kubernetes microservice communication be
modelled and analysed in real time to detect and mitigate sophisticated DDoS attacks
that evade conventional, feature-based detection systems?}
\end{quote}

Specifically, the following limitations of existing solutions motivate this work:

\begin{enumerate}
  \item \textbf{Flat traffic analysis:} Existing ML-based detection systems (e.g., Random
  Forest~\cite{rf_ddos}, XGBoost~\cite{xgboost_ids}) extract features from individual
  flows, discarding the relational context between communicating endpoints.

  \item \textbf{Limited intra-cluster visibility:} Industry solutions such as Cloudflare
  and AWS Shield operate at the network perimeter and cannot observe encrypted, internal
  pod-to-pod traffic within a Kubernetes cluster.

  \item \textbf{Static mitigation policies:} Rate limiting and traffic scrubbing rules are
  typically static and threshold-based, leading to either excessive false positives during
  legitimate traffic spikes or insufficient protection against slow, stealthy attacks.

  \item \textbf{No graph-aware Kubernetes security:} Prior work on graph-based IDS
  (e.g., E-GraphSAGE~\cite{lo2022egraphsage}) applies GNNs to IP-flow graphs rather than
  to Kubernetes service topology graphs, missing the structural context of
  orchestrator-managed communication.
\end{enumerate}

% ─────────────────────────────────────────────
\section{Research Questions}
\label{sec:rqs}

This project is guided by the following research questions:

\begin{description}
  \item[\textbf{RQ1}] How can eBPF-based network observability and dynamic graph modelling
  be combined to accurately represent microservice communication patterns in a Kubernetes
  environment?

  \item[\textbf{RQ2}] How effectively can graph-based anomaly detection --- using
  techniques such as Graph Autoencoders and GraphSAGE --- identify sophisticated DDoS
  attacks, including low-and-slow and application-layer variants, compared with
  conventional traffic-based detection methods?

  \item[\textbf{RQ3}] How can adaptive rate limiting, driven by graph-derived anomaly
  scores, improve DDoS mitigation effectiveness while minimising resource consumption and
  impact on legitimate users?
\end{description}

% ─────────────────────────────────────────────
\section{Research Objectives}
\label{sec:objectives}

To answer the above questions, this project pursues the following objectives:

\begin{description}
  \item[\textbf{RO1}] Develop a Kubernetes-native monitoring framework using eBPF and
  graph modelling to capture and represent microservice communication patterns in real time.

  \item[\textbf{RO2}] Design, implement, and evaluate a graph-based anomaly detection
  model capable of detecting volumetric and sophisticated DDoS attacks in Kubernetes
  environments.

  \item[\textbf{RO3}] Develop and evaluate an adaptive rate-limiting mechanism that
  utilises graph-derived anomaly scores to selectively throttle suspicious traffic while
  preserving legitimate service availability.
\end{description}

% ─────────────────────────────────────────────
\section{Proposed Approach}
\label{sec:approach}

The proposed solution, named \sys{}, is a four-phase pipeline:

\begin{enumerate}
  \item \textbf{Observe:} Deploy eBPF probes via Cilium~\cite{cilium} and export
  real-time network flow telemetry from the Kubernetes data plane using the Hubble
  observability layer~\cite{hubble}.

  \item \textbf{Model:} Construct sliding-window dynamic graphs where nodes represent
  pods and services, and edges encode communication metadata (throughput, latency,
  connection count, error rate).

  \item \textbf{Detect:} Apply a Graph Autoencoder (GAE)~\cite{kipf2016vgae} trained on
  normal communication patterns; high reconstruction error signals anomalous behaviour.
  A GraphSAGE~\cite{hamilton2017graphsage} classifier is subsequently trained to
  discriminate attack sub-types.

  \item \textbf{Mitigate:} Use anomaly scores to drive an adaptive mitigation controller
  that enforces dynamic rate limits via XDP/TC eBPF programs and updates
  \texttt{CiliumNetworkPolicy} resources, without resorting to indiscriminate traffic
  dropping.
\end{enumerate}

% ─────────────────────────────────────────────
\section{Expected Contributions}
\label{sec:contributions}

This project makes the following expected contributions:

\begin{enumerate}
  \item A novel \textbf{eBPF-powered graph construction framework} for real-time
  modelling of Kubernetes microservice communication.

  \item A \textbf{graph-based anomaly detection pipeline} that captures relational and
  temporal attack signatures invisible to flow-based detectors.

  \item An \textbf{adaptive, score-driven mitigation controller} that replaces static
  rate-limiting thresholds with dynamic, per-service enforcement policies.

  \item A \textbf{reproducible experimental evaluation} on a Kubernetes testbed,
  including a labelled dataset of simulated DDoS scenarios against microservice workloads.

  \item An open-source prototype suitable for further academic research and potential
  deployment in cloud-native security pipelines.
\end{enumerate}

% ─────────────────────────────────────────────
\section{Scope and Limitations}
\label{sec:scope}

The project focuses on Linux-based Kubernetes clusters using eBPF-capable kernels
(v5.8+). The experimental evaluation is conducted on a simulated testbed rather than
production infrastructure. Encrypted application-layer payloads are not inspected;
detection relies solely on network flow metadata observable at the kernel level. The
framework is evaluated against selected DDoS attack types and does not claim comprehensive
coverage of all attack vectors.

% ─────────────────────────────────────────────
\section{Report Structure}
\label{sec:structure}

The remainder of this report is organised as follows:

\begin{itemize}
  \item \textbf{Chapter~\ref{ch:background}} provides background on Kubernetes networking,
  eBPF, graph theory fundamentals, and GNN architectures.

  \item \textbf{Chapter~\ref{ch:literature}} presents the literature review, analysing
  related work across eBPF security, DDoS detection, and graph-based anomaly detection.

  \item \textbf{Chapter~\ref{ch:design}} presents the system architecture and first-cut
  design of \sys{}.

  \item \textbf{Chapter~\ref{ch:implementation}} describes the implementation in detail.

  \item \textbf{Chapter~\ref{ch:evaluation}} presents the experimental methodology and
  results.

  \item \textbf{Chapter~\ref{ch:discussion}} discusses findings, limitations, and future
  directions.

  \item \textbf{Chapter~\ref{ch:conclusion}} summarises the project contributions.
\end{itemize}
